فرض کنید $p_1, p_2, \dots, p_n$ اولین $n$ عدد اول باشند. بنابر قضیه‌ی باقیمانده‌ی چینی، عدد صحیح مثبت $x_0$ وجود دارد که در دستگاه همنهشتی‌های
\[
x + k \equiv 0 \pmod{p_k^2}, \qquad 1 \leq k \leq n
\]
صدق می‌کند. (معادلًا، $x \equiv -k \pmod{p_k^2}$ برای هر $k$.)

قرار دهید $x = x_0 + \prod_{k=1}^n p_k^2$. آنگاه برای هر $k = 1, \dots, n$ داریم
\[
x + k \equiv 0 \pmod{p_k^2},
\]
پس $p_k^2$ عدد $x + k$ را می‌شمارد. به‌ویژه، $x + k$ حداقل دو بار بر $p_k$ بخش‌پذیر است. علاوه‌براین، چون می‌توان $x$ را بزرگ‌تر از هر کران ثابت گرفت، می‌توان فرض کرد $x + k > p_k^2$، که در این صورت $x + k$ عامل مربعی دارد و در نتیجه توان خالص یک عدد اول واحد نیست.

بنابراین $n$ عدد صحیح مثبت متوالی $x+1, x+2, \dots, x+n$ هیچ‌کدام توان صحیح یک عدد اول نیستند.
